\NSPage{091}%
\end{lemma}
\begin{proof}
\begin{EESegment}{EE-TGT-P091-001}
First estimate the primitive without differentiating the argument inside it, which shifts rapidly.
Then check the exact radial identity and use Fourier inversion on the torus to estimate the cutoff
remainder.
\end{EESegment}

\begin{EESegment}{EE-TGT-P091-002}
\noindent\textbf{Step 1. Support and coefficient estimates.} Work in a neighborhood with one fixed
common index, as in Lemma~\ref{lem:6.2}, and put
\NSi{NS-F-P091-001}{\mathcal L=v_r\cdot\partial_y} and
\NSi{NS-F-P091-002}{g=R^e f}. Multiplication by \NSi{NS-F-P091-003}{R^e} leaves the input
estimate unchanged on the shell. To estimate \NSi{NS-F-P091-004}{\mathcal I_cg}, set
\NSi{NS-F-P091-005}{U=R^{d_r}} and extend
\NSi{NS-F-P091-006}{F(U)=g(R)/(d_rR^{d_r-1})} smoothly by zero. The shell stays inside one fixed
compact subset of \NSi{NS-F-P091-007}{R>0}, so this change of variables and its inverse have
bounded derivatives of every fixed order. The exact representation on the half-lines is
\end{EESegment}
\NSFormula{NS-F-P091-008}{display}{8.9}
\begin{equation}
\begin{aligned}
A_-(U,y)&=\int_{-\infty}^{0}F(U+u,y+Muv_r)\,du,\\
A_+(U,y)&=\int_{0}^{\infty}F(U+u,y+Muv_r)\,du,\\
\mathcal Jg&=A_-+A_+,\qquad
\mathcal I_cg=(1-\chi_m)A_- -\chi_mA_+.
\end{aligned}
\tag{8.9}\label{eq:8.9}
\end{equation}
\begin{EESegment}{EE-TGT-P091-003}
For fixed \NSi{NS-F-P091-009}{u}, the shift depends on neither
\NSi{NS-F-P091-010}{U} nor \NSi{NS-F-P091-011}{(Z,T)}. So a coefficient derivative acts only
on \NSi{NS-F-P091-012}{F} and \NSi{NS-F-P091-013}{\chi_m}, and brings in no factor of
\NSi{NS-F-P091-014}{M}. The input's extension by zero includes derivatives of the moving shell,
so taking derivatives creates no boundary terms.

Near the inner edge, only the left integral appears. Write
\NSi{NS-F-P091-015}{d_a=\log(X/X_a)} and \NSi{NS-F-P091-016}{d_b=\log(X_b/X)}, so
\NSi{NS-F-P091-017}{\zeta=\exp(-a_a/d_a^2-a_b/d_b^2)} and
\NSi{NS-F-P091-018}{\delta=\min(1,d_a,d_b)}, where
\NSi{NS-F-P091-019}{a_a,a_b} are fixed and positive. For every fixed
\NSi{NS-F-P091-020}{B}, the function
\NSi{NS-F-P091-021}{s\mapsto e^{-a_a/s^2}s^{-B}} is increasing when
\NSi{NS-F-P091-022}{s} is positive and small enough. Along the left integration segment, this
bounds the input weight by the same kind of weight at the endpoint; the factors from the opposite
edge are bounded and comparable. Near the outer edge, use the right integral and the matching
monotonicity for \NSi{NS-F-P091-023}{d_b}. In the interior, the output weight has a positive lower
bound. Each integration runs over a bounded length. The same facts apply after any fixed coefficient
derivative, and prove the \NSi{NS-F-P091-024}{\mathsf M_\alpha} estimate. Formula~\eqref{eq:8.9}
also shows directly that the output is zero below and above the shell. Multiplying by
\NSi{NS-F-P091-025}{R^{\pm e}} costs only fixed constants there.
\end{EESegment}

\begin{EESegment}{EE-TGT-P091-004}
\noindent\textbf{Step 2. The radial equation.} Acting on the integrand in \eqref{eq:8.9}, the
operator \NSi{NS-F-P091-026}{\partial_U+M\mathcal L} is the same as taking its
\NSi{NS-F-P091-027}{u}-derivative. It therefore sends
\NSi{NS-F-P091-028}{A_-} to \NSi{NS-F-P091-029}{F}, sends
\NSi{NS-F-P091-030}{A_+} to \NSi{NS-F-P091-031}{-F}, and sends
\NSi{NS-F-P091-032}{\mathcal Jg} to zero. Multiply by
\NSi{NS-F-P091-033}{d_rR^{d_r-1}}, then use the product rule for
\NSi{NS-F-P091-034}{R^{-e}}. This gives \eqref{eq:8.7}. In physical variables, the same
calculation is \NSi{NS-F-P091-035}{\mathfrak r\mathcal I_cg=g-(\partial_r\chi_m)\mathcal Jg}.
\end{EESegment}

\begin{EESegment}{EE-TGT-P091-005}
\noindent\textbf{Step 3. The cutoff remainder under the moment condition.} The radial identity and
the support bounds are now proved. To estimate the cutoff term still left, use the cancellation in
the full radial integral. Haar measure does not change under translation, so
\NSi{NS-F-P091-036}{\langle\mathcal Jg\rangle_Y=\int\langle g\rangle_Y\,dR}. Under the weighted
moment condition, this is zero for \NSi{NS-F-P091-037}{g=R^ef}. If
\NSi{NS-F-P091-038}{H} is a zero-mean function on the torus, the Diophantine inequality
\eqref{eq:6.7} says that the absolute value of either dot product is at least one fixed positive
constant divided by one plus the size of the nonzero integer frequency. It gives
\end{EESegment}
\NSFormula{NS-F-P091-039}{display}{8.10}
\begin{equation}
\|\mathcal L^{-p}H\|_{C_y^m}\le C_{m,p}\|H\|_{C_y^{m+p+3}},
\qquad p\ge1.
\tag{8.10}\label{eq:8.10}
\end{equation}
\begin{EESegment}{EE-TGT-P091-006}
The Fourier multiplier here is
\NSi{NS-F-P091-040}{(2\pi i v_r\cdot k)^{-p}}. After
\NSi{NS-F-P091-041}{m+p+3} derivatives of \NSi{NS-F-P091-042}{H} have been estimated, the
series still to be summed is bounded by
\NSi{NS-F-P091-043}{\sum_{k\in\mathbb Z^2}(1+|k|)^{-3}}. Slow and radial derivatives commute
with this multiplier. The zero-average part of the full-line integral therefore obeys the exact
identity
\end{EESegment}
\NSFormula{NS-F-P091-044}{display}{none}
\[
\mathcal Jg=(-M^{-1})^p\int_{\mathbb R}\partial_U^p\mathcal L^{-p}F^\circ
   (U+u,y+Muv_r)\,du.
\]
\begin{EESegment}{EE-TGT-P091-007}
This follows by integrating the \NSi{NS-F-P091-045}{u}-derivative of
\NSi{NS-F-P091-046}{\mathcal L^{-1}F^\circ}, then repeating the same step; compact support makes
every endpoint term zero. The zero Fourier mode left out of this formula has a full-line
\NSPage{092}%
integral that is identically zero at every slow point. Averaging, integration, and slow
differentiation of the coefficients commute on the smooth extensions by zero, so every derivative
of that mode also has zero integral. The same fixed-shift argument handles any further coefficient
derivatives. For the output derivative indexed by \NSi{NS-F-P092-001}{I}, it is enough to use
radial input derivatives through \NSi{NS-F-P092-002}{I_R+p}, torus derivatives through
\NSi{NS-F-P092-003}{|I_y|+p+3}, and the slow derivatives named by
\NSi{NS-F-P092-004}{I}. Since
\NSi{NS-F-P092-005}{M^{-1}\le C\varepsilon^{\kappa_s}S_*^{\rho_g}} and
\NSi{NS-F-P092-006}{\partial_R\chi_m} is supported in the interior, this proves
\eqref{eq:8.8}. Fix a physical normalization and a physical derivative order, and let
\NSi{NS-F-P092-007}{L\ge0} bound the resulting loss of powers of
\NSi{NS-F-P092-008}{q}, including the loss from differentiating after evaluation at the phase. The
graph operators give a finite \NSi{NS-F-P092-009}{L} of this kind that does not depend on
\NSi{NS-F-P092-010}{p}. Given any flatness order \NSi{NS-F-P092-011}{N\ge0}, choose
\NSi{NS-F-P092-012}{p} so that
\NSi{NS-F-P092-013}{h(\alpha+p\kappa_s)>N+L}. Since
\NSi{NS-F-P092-014}{\varepsilon=Q^h} and \NSi{NS-F-P092-015}{q\asymp Q}, this strict excess
absorbs the fixed powers of \NSi{NS-F-P092-016}{S_*} and gives an
\NSi{NS-F-P092-017}{O(q^N)} bound for that physical derivative. No estimate uniform in
\NSi{NS-F-P092-018}{p} is being claimed. Finally, when \NSi{NS-F-P092-019}{f} does not depend
on \NSi{NS-F-P092-020}{Y},
\NSi{NS-F-P092-021}{\mathcal J(R^ef)=\int R^ef\,dR}, which is zero under the stated weighted
integral condition.
\end{EESegment}
\end{proof}

\begin{EESegment}{EE-TGT-P092-001}
Whenever the weighted mean integral is zero, the inverse is exact apart from a flat cutoff
remainder. This flatness uses estimates at every fixed derivative order. Frequencies comparable to
\NSi{NS-F-P092-022}{M} can come close to resonance with \NSi{NS-F-P092-023}{v_r}, meaning that
the dot product with the frequency can be small and division by it can be large. The extra
derivatives in the torus variables in \eqref{eq:8.10} control their small Fourier coefficients. A
fixed finite regularity bound would give only a fixed finite order of flatness.
\end{EESegment}

\begin{EESegment}{EE-TGT-P092-002}
\subsection{Reconstructing the pressure and making divergence-free velocity corrections.}\label{sec:8.3}
This subsection applies the radial inverse from Lemma~\ref{lem:8.2} first to the pressure
\NSi{NS-F-P092-024}{p_m}, then to an azimuthal vector potential that produces a divergence-free
velocity increment; see Proposition~\ref{prop:8.3}. The axial component of that increment equals
the prescribed source \NSi{NS-F-P092-025}{\gamma_d}, apart from the cutoff remainder in
\eqref{eq:8.14}. The pressure source may have a nonzero radial integral after its auxiliary average is
taken. Subtracting a fixed bump times that integral separates out the obstruction, and the source
left over can then be integrated with compact support, as in \eqref{eq:8.12}.
\end{EESegment}

\begin{EESegment}{EE-TGT-P092-003}
Let \NSi{NS-F-P092-026}{\mathcal I_{\mathrm{mean}}} be the interval in
\NSi{NS-F-P092-027}{X} reserved by Theorem~\ref{thm:4.6}(vi), and set
\end{EESegment}
\NSFormula{NS-F-P092-028}{display}{none}
\[
I_m=\{\sqrt{2X}:X\in \mathcal I_{\mathrm{mean}}\}.
\]
\begin{EESegment}{EE-TGT-P092-004}
The \emph{mean patch} is the region \NSi{NS-F-P092-029}{r/\sqrt q\in I_m}. Choose a fixed bump
\NSi{NS-F-P092-030}{\widehat\rho\in C_c^\infty(I_m)} with
\NSi{NS-F-P092-031}{\int\widehat\rho(x)\,dx=1}, and define
\end{EESegment}
\NSFormula{NS-F-P092-032}{display}{none}
\[
\rho_{\mathrm{phys}}(r,z,t)=q^{-1/2}\widehat\rho(r/\sqrt q),
\qquad \rho=Q^{1/2}\rho_{\mathrm{phys}}.
\]
\begin{EESegment}{EE-TGT-P092-005}
This bump does not depend on the angular or auxiliary variables, and
\NSi{NS-F-P092-033}{\int\rho_{\mathrm{phys}}\,dr=\int\rho\,dR=1}. It is fixed independently of
the velocity corrections.
\end{EESegment}

\begin{EESegment}{EE-TGT-P092-006}
The operator \NSi{NS-F-P092-034}{\mathsf T_0} reconstructs the pressure, while
\NSi{NS-F-P092-035}{\mathsf T_1} reconstructs an azimuthal vector potential. The operators
\NSi{NS-F-P092-036}{\mathsf T_1} and \NSi{NS-F-P092-037}{\mathsf T_2} also recover the radial
fluxes of axial and azimuthal momentum. Write \NSi{NS-F-P092-038}{\Psi} for the azimuthal
coefficient of the vector potential. With a star marking a normalized representative, the physical
normalizations are
\end{EESegment}
\NSFormula{NS-F-P092-039}{display}{8.11}
\begin{equation}
\begin{gathered}
u_*=Q^A u_{\mathrm{phys}},\qquad
g_*=Q^{2A+1/2}g_{\mathrm{phys}},\qquad
p_*=Q^{2A}p_{\mathrm{phys}},\\
\Psi_*=Q^{A-1/2}\Psi_{\mathrm{phys}},\qquad
\rho_*=Q^{1/2}\rho_{\mathrm{phys}}.
\end{gathered}
\tag{8.11}\label{eq:8.11}
\end{equation}
\begin{EESegment}{EE-TGT-P092-007}
In particular,
\NSi{NS-F-P092-040}{Q^{2A}\mathcal I_{c,\mathrm{phys}}g_{\mathrm{phys}}=\mathcal I_cg_*}.
When the \NSi{NS-F-P092-041}{\varepsilon} exponent is said to be preserved, it refers to these
normalized representatives and already includes the length of integration.
\end{EESegment}

\begin{proposition}\label{prop:8.3}
\begin{EESegment}{EE-TGT-P092-008}
Fix the radial operators from Lemma~\ref{lem:8.2} and the bump
\NSi{NS-F-P092-042}{\rho} defined above.
\NSPage{093}%
\end{EESegment}
\begin{enumerate}[label=(\roman*)]
\item
\begin{EESegment}{EE-TGT-P093-001}
\emph{Pressure.} For a smooth shell-supported scalar source
\NSi{NS-F-P093-001}{g_r(R,Z,T,y)}, define the radial defect
\NSi{NS-F-P093-002}{\mathcal P} and the pressure \NSi{NS-F-P093-003}{p_m} by
\end{EESegment}
\NSFormula{NS-F-P093-004}{display}{8.12}
\begin{equation}
\mathcal P=\int\langle g_r\rangle_Y\,dR,
\qquad p_m=\mathsf T_0(g_r-\rho\mathcal P),
\qquad \int\rho\,dR=1.
\tag{8.12}\label{eq:8.12}
\end{equation}
\begin{EESegment}{EE-TGT-P093-002}
Then
\end{EESegment}
\NSFormula{NS-F-P093-005}{display}{8.13}
\begin{equation}
D_rp_m-g_r=-\rho\mathcal P-\mathcal A_0(g_r-\rho\mathcal P),
\qquad \partial_R\langle p_m\rangle_Y=\langle g_r\rangle_Y-\rho\mathcal P.
\tag{8.13}\label{eq:8.13}
\end{equation}
\begin{EESegment}{EE-TGT-P093-003}
If \NSi{NS-F-P093-006}{g_r\in \mathsf M_\alpha}, then
\NSi{NS-F-P093-007}{\mathcal P\in \mathsf S_\alpha} and
\NSi{NS-F-P093-008}{p_m\in \mathsf M_\alpha}. Changing \NSi{NS-F-P093-009}{g_r} by a term
in \NSi{NS-F-P093-010}{\mathsf M_\alpha} changes the pressure by a term in the same class. Under
this class assumption, the cutoff remainder in \eqref{eq:8.13} has zero auxiliary mean and is flat
as \NSi{NS-F-P093-011}{q\downarrow0} at every fixed derivative order.
\end{EESegment}

\item
\begin{EESegment}{EE-TGT-P093-004}
\emph{Axial velocity.} Let \NSi{NS-F-P093-012}{\gamma_d(R,Z,T,y)} be a smooth, shell-supported desired axial increment
such that \NSi{NS-F-P093-013}{\int R\langle\gamma_d\rangle_Y\,dR=0} at every
\NSi{NS-F-P093-014}{(Z,T)}. Using \eqref{eq:8.4}, define the physical azimuthal coefficient of the
vector potential by
\NSi{NS-F-P093-015}{\Psi=-r^{-1}\mathcal I_c(r\gamma_{d,\mathrm{phys}})}. Define its normalized
velocity increments by
\end{EESegment}
\NSFormula{NS-F-P093-016}{display}{8.14}
\begin{equation}
\Psi_*=\mathsf T_1\gamma_d,
\qquad \Delta\beta=-\varepsilon\partial_Z\Psi_*,
\qquad \Delta\gamma=(D_r+1/R)\Psi_*=\gamma_d-\mathcal A_1\gamma_d.
\tag{8.14}\label{eq:8.14}
\end{equation}
\begin{EESegment}{EE-TGT-P093-005}
The actual increment \NSi{NS-F-P093-017}{(\Delta\beta,0,\Delta\gamma)} is divergence-free and
satisfies \NSi{NS-F-P093-018}{\langle\Delta\gamma\rangle_Y=\langle\gamma_d\rangle_Y}. If
\NSi{NS-F-P093-019}{\gamma_d\in \mathsf M_\alpha}, then
\NSi{NS-F-P093-020}{\Psi_*,\Delta\gamma\in \mathsf M_\alpha} and
\NSi{NS-F-P093-021}{\Delta\beta\in \mathsf M_{\alpha+1}}. If
\NSi{NS-F-P093-022}{\gamma_d} does not depend on \NSi{NS-F-P093-023}{Y}, its zero weighted
integral makes \NSi{NS-F-P093-024}{\Delta\gamma=\gamma_d} at every point. The pressure and vector
potential are compactly supported inside the active shell, and their support in
\NSi{NS-F-P093-025}{(Z,T)} stays inside the projection of the source support.
\end{EESegment}
\end{enumerate}
\end{proposition}

\begin{proof}
\begin{EESegment}{EE-TGT-P093-006}
\noindent\textbf{Part (i). Pressure.} The source in \eqref{eq:8.12} has zero radial integral after
its auxiliary mean is taken, since
\end{EESegment}
\NSFormula{NS-F-P093-026}{display}{none}
\[
\int\langle g_r-\rho\mathcal P\rangle_Y\,dR
=\mathcal P-\mathcal P\int\rho\,dR=0.
\]
\begin{EESegment}{EE-TGT-P093-007}
Apply Lemma~\ref{lem:8.2} with \NSi{NS-F-P093-027}{e=0}, then take the auxiliary average. The
moment \NSi{NS-F-P093-028}{\mathcal P} belongs to \NSi{NS-F-P093-029}{\mathsf S_\alpha} when
\NSi{NS-F-P093-030}{g_r\in \mathsf M_\alpha}. To see this,
integrate its derivatives of any fixed order over the bounded shell, using the smooth extension by
zero and the fact that \NSi{NS-F-P093-031}{\zeta\delta^{-B}} is bounded for every finite
\NSi{NS-F-P093-032}{B}. The interior bump \NSi{NS-F-P093-033}{\rho\mathcal P} belongs to
\NSi{NS-F-P093-034}{\mathsf M_\alpha}. Since \NSi{NS-F-P093-035}{\rho} is fixed independently
of the current velocity, the same argument proves the statement about changes and completes part
(i).
\end{EESegment}

\begin{EESegment}{EE-TGT-P093-008}
\noindent\textbf{Part (ii). Axial velocity.} The azimuthal potential produces the two physical
velocity components \NSi{NS-F-P093-036}{-\partial_z\Psi} and
\NSi{NS-F-P093-037}{(\mathfrak r+1/r)\Psi}. The operators
\NSi{NS-F-P093-038}{\mathfrak r} and \NSi{NS-F-P093-039}{\partial_z} commute when they
differentiate after evaluation at the phase map, so the divergence of these two components is
exactly zero:
\end{EESegment}
\NSFormula{NS-F-P093-040}{display}{none}
\[
(\mathfrak r+r^{-1})(-\partial_z\Psi)+\partial_z\bigl((\mathfrak r+r^{-1})\Psi\bigr)=0.
\]
\begin{EESegment}{EE-TGT-P093-009}
The weighted primitive identity gives the last equality in \eqref{eq:8.14}, including its sign. The
zero weighted axial-flux integral makes the auxiliary average of its cutoff remainder zero. Finally,
\NSi{NS-F-P093-041}{Q^{1/2-D}=Q^h=\varepsilon}, which supplies the extra factor
\NSi{NS-F-P093-042}{\varepsilon} in \NSi{NS-F-P093-043}{\Delta\beta} in \eqref{eq:8.14}. The
remaining claims follow from Lemma~\ref{lem:8.2}, and part (ii) is complete. No slow derivative is
moved through \NSi{NS-F-P093-044}{\mathcal I_c} in this argument.
\end{EESegment}
\end{proof}

\begin{EESegment}{EE-TGT-P093-010}
\subsection{The three compatibility defects in the integrated equations.}\label{sec:8.4}
This subsection finds the integral defects left by the pressure reconstruction in \eqref{eq:8.13}
and by the two tangential equations in \eqref{eq:8.3}. These defects are the quantities that must be
removed before the required compactly supported corrections can exist. Tangential stresses are then
built so that their radial divergences cancel the auxiliary-averaged residuals apart from two
explicit bump terms; see Corollary~\ref{cor:8.5}. The pressure reconstruction leaves the integral
\NSi{NS-F-P093-045}{\mathcal P} from \eqref{eq:8.12} to be corrected. In the same way, the
auxiliary averages of the tangential residuals must have zero weighted radial integrals before
compactly supported stresses can represent them, by Lemma~\ref{lem:8.2}. The moment constraints
in \eqref{eq:8.2} remove the time terms and the axial-viscosity terms from these integrals. The axial
fluxes left after that give two more defects, as Equations~\eqref{eq:8.15} and \eqref{eq:8.16} show.
\end{EESegment}

\NSPage{094}%
\begin{EESegment}{EE-TGT-P094-001}
Along with \NSi{NS-F-P094-001}{\mathcal P} from \eqref{eq:8.12}, define the two flux defects and
the second radial moment \NSi{NS-F-P094-002}{c_\rho} of \NSi{NS-F-P094-003}{\rho} by
\end{EESegment}
\NSFormula{NS-F-P094-004}{display}{8.15}
\begin{equation}
\begin{aligned}
\mathcal J_\theta
  &=\int R^2\langle Gv+V\gamma+\gamma v+W_{z\theta}\rangle_Y\,dR,\\
\mathcal J_z
  &=\int R\langle2G\gamma+\gamma^2+W_{zz}\rangle_Y\,dR
    -\frac12\int R^2\langle g_r\rangle_Y\,dR,
&c_\rho&=\frac12\int R^2\rho\,dR.
\end{aligned}
\tag{8.15}\label{eq:8.15}
\end{equation}

\begin{proposition}\label{prop:8.4}
\begin{EESegment}{EE-TGT-P094-002}
Let \NSi{NS-F-P094-005}{(\beta,v,\gamma)} and \NSi{NS-F-P094-006}{w} be the smooth corrections
from \eqref{eq:8.1}, supported in the shell. Suppose that the two moments in \eqref{eq:8.2} are zero
at every \NSi{NS-F-P094-007}{(Z,T)}, and reconstruct \NSi{NS-F-P094-008}{p_m} from their radial
source \NSi{NS-F-P094-009}{g_r} by \eqref{eq:8.12}. For the residual components
\NSi{NS-F-P094-010}{E_\theta,E_z} from \eqref{eq:8.3} and the defects
\NSi{NS-F-P094-011}{\mathcal J_\theta,\mathcal J_z,c_\rho} from \eqref{eq:8.15},
\end{EESegment}
\NSFormula{NS-F-P094-012}{display}{8.16}
\begin{equation}
\int R^2\langle E_\theta\rangle_Y\,dR
=\varepsilon\partial_Z\mathcal J_\theta,
\qquad
\int R\langle E_z\rangle_Y\,dR
=\varepsilon\partial_Z(\mathcal J_z+c_\rho\mathcal P).
\tag{8.16}\label{eq:8.16}
\end{equation}
\end{proposition}

\begin{proof}
\begin{EESegment}{EE-TGT-P094-003}
After auxiliary averaging, \NSi{NS-F-P094-013}{D_r} becomes
\NSi{NS-F-P094-014}{\partial_R}, and \NSi{NS-F-P094-015}{\mathfrak t_*} becomes
\NSi{NS-F-P094-016}{-\varepsilon\partial_T}. Multiply each radial divergence in
\eqref{eq:8.3}, including the one from the residual stress tensor, by its stated weight and integrate.
Each one becomes a term at the endpoints, and compact support makes those terms zero. Two
integrations by parts make the radial viscous integrals zero.
\end{EESegment}
\NSFormula{NS-F-P094-017}{display}{none}
\[
\int R^2(v_{RR}+R^{-1}v_R-R^{-2}v)\,dR
=(2-1-1)\int v\,dR=0,
\]
\NSFormula{NS-F-P094-018}{display}{none}
\[
\int R(\gamma_{RR}+R^{-1}\gamma_R)\,dR
=-\int\gamma_R\,dR+\int\gamma_R\,dR=0.
\]
\begin{EESegment}{EE-TGT-P094-004}
These formulas are applied to the auxiliary averages. The time terms and axial-viscosity terms are
derivatives of the two integrals that are already zero. The compactly supported pressure
reconstruction fixes the pressure moment as follows.
\end{EESegment}
\NSFormula{NS-F-P094-019}{display}{none}
\[
\int R\langle p_m\rangle_Y\,dR
=-\frac12\int R^2\partial_R\langle p_m\rangle_Y\,dR
=-\frac12\int R^2\langle g_r\rangle_Y\,dR+c_\rho\mathcal P.
\]
\begin{EESegment}{EE-TGT-P094-005}
Putting this expression into the axial fluxes that remain gives \eqref{eq:8.16}. Since
\NSi{NS-F-P094-020}{\rho} is scaled using the moving \NSi{NS-F-P094-021}{q},
\NSi{NS-F-P094-022}{c_\rho} may depend on \NSi{NS-F-P094-023}{(Z,T)}. For this reason, its full
product with \NSi{NS-F-P094-024}{\mathcal P} stays inside
\NSi{NS-F-P094-025}{\partial_Z}.
\end{EESegment}
\end{proof}

\begin{EESegment}{EE-TGT-P094-006}
The moments in physical variables and their normalized representatives are related by
\end{EESegment}
\NSFormula{NS-F-P094-026}{display}{8.17}
\begin{equation}
\begin{gathered}
(M_\theta)_*=Q^{A-3/2}(M_\theta)_{\mathrm{phys}},
\qquad (M_z)_*=Q^{A-1}(M_z)_{\mathrm{phys}},\\
\mathcal P_*=Q^{2A}\mathcal P_{\mathrm{phys}},
\qquad (c_\rho)_*=Q^{-1}(c_\rho)_{\mathrm{phys}},\\
(\mathcal J_\theta)_*=Q^{2A-3/2}(\mathcal J_\theta)_{\mathrm{phys}},
\qquad (\mathcal J_z)_*=Q^{2A-1}(\mathcal J_z)_{\mathrm{phys}}.
\end{gathered}
\tag{8.17}\label{eq:8.17}
\end{equation}
\begin{EESegment}{EE-TGT-P094-007}
The constraints that set an integral to zero do not change under normalization. Every nonzero
correction target uses the normalization for its own row.
\end{EESegment}

\begin{EESegment}{EE-TGT-P094-008}
The two tangential residual moments are therefore axial derivatives of
\NSi{NS-F-P094-027}{\mathcal J_\theta} and
\NSi{NS-F-P094-028}{\mathcal J_z+c_\rho\mathcal P}. Subtract a fixed normalized bump times each
weighted residual integral, then apply the radial inverse. These sources do not depend on the
auxiliary variable, and their weighted integrals are zero, so the inverse is exact. The resulting pair
will be the target for covariance added by the waves. The weighted radial divergence of that
covariance enters \eqref{eq:8.3} with a positive sign through
\NSi{NS-F-P094-029}{W_{r\theta}} and \NSi{NS-F-P094-030}{W_{rz}}. Fix the interior profiles
\NSi{NS-F-P094-031}{\widehat\sigma_\theta,\widehat\sigma_z} so that
\NSi{NS-F-P094-032}{\int\xi^2\widehat\sigma_\theta(\xi)\,d\xi
=\int\xi\widehat\sigma_z(\xi)\,d\xi=1}, and define
\end{EESegment}
\NSFormula{NS-F-P094-033}{display}{none}
\[
\sigma_{\theta,\mathrm{phys}}=q^{-3/2}\widehat\sigma_\theta(r/\sqrt q),
\qquad
\sigma_{z,\mathrm{phys}}=q^{-1}\widehat\sigma_z(r/\sqrt q).
\]
\begin{EESegment}{EE-TGT-P094-009}
Their supports lie inside the mean patch. In each chart, use
\NSi{NS-F-P094-034}{\sigma_\theta=Q^{3/2}\sigma_{\theta,\mathrm{phys}}} and
\NSi{NS-F-P094-035}{\sigma_z=Q\sigma_{z,\mathrm{phys}}}, so
\NSi{NS-F-P094-036}{\int R^2\sigma_\theta\,dR=\int R\sigma_z\,dR=1}. This gives the same
physical bump functions wherever charts overlap.
\end{EESegment}

\NSPage{095}%
\begin{corollary}\label{cor:8.5}
\begin{EESegment}{EE-TGT-P095-001}
Under the assumptions of Proposition~\ref{prop:8.4}, define the covariance targets
\end{EESegment}
\NSFormula{NS-F-P095-001}{display}{8.18}
\begin{equation}
\begin{aligned}
H_\theta&=-\mathsf T_2\bigl(\langle E_\theta\rangle_Y
  -\sigma_\theta\varepsilon\partial_Z\mathcal J_\theta\bigr),\\
H_z&=-\mathsf T_1\bigl(\langle E_z\rangle_Y
  -\sigma_z\varepsilon\partial_Z(\mathcal J_z+c_\rho\mathcal P)\bigr).
\end{aligned}
\tag{8.18}\label{eq:8.18}
\end{equation}
\begin{EESegment}{EE-TGT-P095-002}
They do not depend on the auxiliary variable, are compactly supported inside the active shell, and
satisfy
\end{EESegment}
\NSFormula{NS-F-P095-002}{display}{none}
\begin{align*}
(D_r+2/R)H_\theta&=-\langle E_\theta\rangle_Y
  +\sigma_\theta\varepsilon\partial_Z\mathcal J_\theta,\\
(D_r+1/R)H_z&=-\langle E_z\rangle_Y
  +\sigma_z\varepsilon\partial_Z(\mathcal J_z+c_\rho\mathcal P).
\end{align*}
\begin{EESegment}{EE-TGT-P095-003}
The radial primitives act while \NSi{NS-F-P095-003}{(Z,T)} is held fixed. It follows that
\NSi{NS-F-P095-004}{H_\theta} and \NSi{NS-F-P095-005}{H_z} are zero whenever their matching
integrands in \eqref{eq:8.18} are zero for every \NSi{NS-F-P095-006}{R}. The construction does
not enlarge the support in the axial variable \NSi{NS-F-P095-007}{Z}.
\end{EESegment}
\end{corollary}

\begin{proof}
\begin{EESegment}{EE-TGT-P095-004}
The integrated identities and the normalization of the two bumps give
\end{EESegment}
\NSFormula{NS-F-P095-008}{display}{none}
\begin{align*}
\int R^2\bigl(\langle E_\theta\rangle_Y
  -\sigma_\theta\varepsilon\partial_Z\mathcal J_\theta\bigr)\,dR&=0,\\
\int R\bigl(\langle E_z\rangle_Y
  -\sigma_z\varepsilon\partial_Z(\mathcal J_z+c_\rho\mathcal P)\bigr)\,dR&=0.
\end{align*}
\begin{EESegment}{EE-TGT-P095-005}
Both sources are independent of the auxiliary variable. Lemma~\ref{lem:8.2} makes their cutoff
remainders zero, which proves the two identities and the statements about support.
\end{EESegment}
\end{proof}

\begin{EESegment}{EE-TGT-P095-006}
These radial-divergence identities leave the factors
\NSi{NS-F-P095-009}{\varepsilon\partial_Z\mathcal J_\theta} and
\NSi{NS-F-P095-010}{\varepsilon\partial_Z(\mathcal J_z+c_\rho\mathcal P)} in the bump terms
that remain. The three defects can therefore be corrected at fixed
\NSi{NS-F-P095-011}{(Z,T)} without losing the displayed factors of
\NSi{NS-F-P095-012}{\varepsilon\partial_Z}.
\end{EESegment}

\begin{EESegment}{EE-TGT-P095-007}
\subsection{Inverting the auxiliary-time derivative on zero-mean functions.}\label{sec:8.5}
Corollary~\ref{cor:8.5} reduces the auxiliary-averaged correction to the three integral defects and
a compactly supported stress. The part of the mean residual left after its auxiliary average is
removed, \NSi{NS-F-P095-013}{(E_\theta^\circ,E_z^\circ)}, has zero auxiliary average at every
slow point by \eqref{eq:3.9}. Its correction comes from the fast time derivative. The Fourier
multiplier for that derivative is invertible at every nonzero torus frequency by \eqref{eq:6.7}. The
axial velocity increment is then made from its azimuthal potential through \eqref{eq:8.14}, which
keeps the velocity incompressible.
\end{EESegment}

\begin{EESegment}{EE-TGT-P095-008}
Recall from Equations~\eqref{eq:6.4} and \eqref{eq:6.6} that
\NSi{NS-F-P095-014}{N_{\mathrm{abs}}=v_t\cdot\partial_Y} acts on the original auxiliary
coordinate. On a common torus \NSi{NS-F-P095-015}{y=Y_{i_0}}, the normalized fast time
derivative is \NSi{NS-F-P095-016}{c_{i_0}N_{i_0}}, where
\NSi{NS-F-P095-017}{N_{i_0}=v_t\cdot\partial_y} and
\NSi{NS-F-P095-018}{c_{i_0}=T_g^{i_0}Q^{1+h}}. The normalized physical time derivative splits
as \NSi{NS-F-P095-019}{\mathfrak t_*=-\varepsilon\partial_T+c_{i_0}N_{i_0}}. Its slow part
\NSi{NS-F-P095-020}{-\varepsilon\partial_T} differentiates the coefficient's explicit dependence
on time while \NSi{NS-F-P095-021}{y} stays fixed. Its fast part
\NSi{NS-F-P095-022}{c_{i_0}N_{i_0}} differentiates the auxiliary oscillations along the physical
phase map. The fast part is inverted here, while the slow part stays in the residual.
\end{EESegment}

\begin{lemma}\label{lem:8.6}
\begin{EESegment}{EE-TGT-P095-009}
Let \NSi{NS-F-P095-023}{N=v_t\cdot\partial_y} on
\NSi{NS-F-P095-024}{\mathbb T^2}, with \NSi{NS-F-P095-025}{v_t} as in \eqref{eq:6.2}, and let
\NSi{NS-F-P095-026}{F} be a smooth function with zero Haar mean. The equation
\NSi{NS-F-P095-027}{N\varphi=F} has exactly one smooth solution with zero mean, namely
\NSi{NS-F-P095-028}{\varphi=N^{-1}F}. For every integer
\NSi{NS-F-P095-029}{m\ge0}, this solution is given by the Fourier series below and satisfies the
stated norm bound.
\end{EESegment}
\NSFormula{NS-F-P095-030}{display}{8.19}
\begin{equation}
N^{-1}F=\sum_{k\in\mathbb Z^2\setminus\{0\}}
\frac{\widehat F(k)}{2\pi i v_t\cdot k}e^{2\pi i k\cdot y},
\qquad
\|N^{-1}F\|_{C_y^m}\le C_m\|F\|_{C_y^{m+4}}.
\tag{8.19}\label{eq:8.19}
\end{equation}
\begin{EESegment}{EE-TGT-P095-010}
For a pair \NSi{NS-F-P095-031}{E_\theta,E_z\in \mathsf M_\alpha} of normalized residual
coefficients that do not depend on the physical angle, use this inverse in the original auxiliary
coordinate \NSi{NS-F-P095-032}{Y}. It defines the physical tangential update
\NSi{NS-F-P095-033}{-N_{\mathrm{abs}}^{-1}
  (E_{\theta,\mathrm{phys}}^\circ,E_{z,\mathrm{phys}}^\circ)}.
\NSPage{096}%
The residuals and velocities here have the normalizations in \eqref{eq:8.11}. Defining the update
in physical coordinates makes its values agree wherever charts overlap. Its chart representatives
are
\end{EESegment}
\NSFormula{NS-F-P096-001}{display}{8.20}
\begin{equation}
\Delta v=-c_{i_0}^{-1}N_{i_0}^{-1}E_\theta^\circ,
\qquad
\gamma_d=-c_{i_0}^{-1}N_{i_0}^{-1}E_z^\circ,
\qquad
c_{i_0}^{-1}\le CS_*.
\tag{8.20}\label{eq:8.20}
\end{equation}
\begin{EESegment}{EE-TGT-P096-001}
Both desired increments \NSi{NS-F-P096-002}{\Delta v,\gamma_d} belong to
\NSi{NS-F-P096-003}{\mathsf M_\alpha}, and each has zero auxiliary mean at every slow and radial
point. Realizing \NSi{NS-F-P096-004}{\gamma_d} through \eqref{eq:8.14} gives the actual velocity
increment
\end{EESegment}
\NSFormula{NS-F-P096-005}{display}{8.21}
\begin{equation}
\begin{gathered}
(\Delta\beta,\Delta v,\Delta\gamma)
=(-\varepsilon\partial_Z\mathsf T_1\gamma_d,\Delta v,\gamma_d-\mathcal A_1\gamma_d),\\
(\Delta\beta,\Delta v,\Delta\gamma)
\in \mathsf M_{\alpha+1}\times \mathsf M_\alpha\times \mathsf M_\alpha.
\end{gathered}
\tag{8.21}\label{eq:8.21}
\end{equation}
\begin{EESegment}{EE-TGT-P096-002}
This field is divergence-free, and adding it leaves both integral constraints in \eqref{eq:8.2}
unchanged. Let \NSi{NS-F-P096-006}{a=-\mathcal A_1\gamma_d} be the cutoff remainder in the
axial increment. Its cancellation equations for the fast-time derivative are
\end{EESegment}
\NSFormula{NS-F-P096-007}{display}{8.22}
\begin{equation}
c_{i_0}N_{i_0}\Delta v=-E_\theta^\circ,
\qquad
c_{i_0}N_{i_0}\Delta\gamma=-E_z^\circ+c_{i_0}N_{i_0}a.
\tag{8.22}\label{eq:8.22}
\end{equation}
\begin{EESegment}{EE-TGT-P096-003}
The angular cancellation is exact. The term
\NSi{NS-F-P096-008}{c_{i_0}N_{i_0}a} is flat at every fixed order of coefficient derivatives.
\end{EESegment}
\end{lemma}

\begin{proof}
\begin{EESegment}{EE-TGT-P096-004}
The divisor bound in \eqref{eq:6.7} gives
\NSi{NS-F-P096-009}{|v_t\cdot k|^{-1}\le C(1+|k|)}. Taking four more derivatives of
\NSi{NS-F-P096-010}{F} than the number requested for the output leaves a summable
\NSi{NS-F-P096-011}{(1+|k|)^{-3}} bound for the two-dimensional Fourier series. This proves
\eqref{eq:8.19}, smoothness, and uniqueness once the zero Fourier coefficient is fixed. The
multiplier commutes with every slow coefficient derivative and preserves real-valued functions. It
acts while \NSi{NS-F-P096-012}{(R,Z,T)} is held fixed, so it preserves slow and radial support and
their smooth extensions by zero, though it does not have to preserve support on the auxiliary torus.
\end{EESegment}

\begin{EESegment}{EE-TGT-P096-005}
Let \NSi{NS-F-P096-013}{P_iF(Y)=F(J_g^iY)}. A mode \NSi{NS-F-P096-014}{k} on the
\NSi{NS-F-P096-015}{i}-torus has frequency \NSi{NS-F-P096-016}{(J_g^i)^Tk} in the original
auxiliary coordinate \NSi{NS-F-P096-017}{Y}. Since
\NSi{NS-F-P096-018}{J_gv_t=T_gv_t},
\end{EESegment}
\NSFormula{NS-F-P096-019}{display}{8.23}
\begin{equation}
N_{\mathrm{abs}}P_i=T_g^iP_iN_i,
\qquad
N_{\mathrm{abs}}^{-1}P_i=T_g^{-i}P_iN_i^{-1}.
\tag{8.23}\label{eq:8.23}
\end{equation}
\begin{EESegment}{EE-TGT-P096-006}
The finite covering preserves the Haar average. The same fact also follows from its injective map
between the frequency lattices. Applying \eqref{eq:8.19} in the common frequency coordinate
\NSi{NS-F-P096-020}{k} gives the same fixed loss of derivatives in every chart.
\end{EESegment}

\begin{EESegment}{EE-TGT-P096-007}
In physical variables, the update is
\NSi{NS-F-P096-021}{-N_{\mathrm{abs}}^{-1}E_{\mathrm{phys}}^\circ}. Its value in chart coordinates
has the factor
\NSi{NS-F-P096-022}{Q^{A-(2A+1/2)}T_g^{-i_0}=Q^{-1-h}T_g^{-i_0}=c_{i_0}^{-1}}, which proves
\eqref{eq:8.20}. The choice of covering level and its bounded difference from a band level give the
bound \NSi{NS-F-P096-023}{CS_*} for this factor. For each output derivative, four additional
derivatives in the torus variables are enough, so the class estimates at every fixed derivative order
keep their \NSi{NS-F-P096-024}{\varepsilon} exponent. Applying the slow-time term
\NSi{NS-F-P096-025}{-\varepsilon\partial_T} to these increments supplies one more factor of
\NSi{NS-F-P096-026}{\varepsilon}.
\end{EESegment}

\begin{EESegment}{EE-TGT-P096-008}
The desired increments have zero auxiliary mean at every radial point. This gives both desired
integral constraints, and making the increment from the azimuthal potential keeps them exactly. The
bounds for \NSi{NS-F-P096-027}{\Delta\beta} and
\NSi{NS-F-P096-028}{\Delta\gamma} in \eqref{eq:8.21} follow from
Proposition~\ref{prop:8.3}(ii). The fast derivative commutes with
\NSi{NS-F-P096-029}{\mathcal I}, \NSi{NS-F-P096-030}{\mathcal J}, and
\NSi{NS-F-P096-031}{\mathcal I_c}, because the shifts are translations and
\NSi{NS-F-P096-032}{\chi_m} does not depend on the torus variable. Apply that derivative to
\NSi{NS-F-P096-033}{\Delta\gamma=\gamma_d+a} to get \eqref{eq:8.22}. Its last term is flat at
every fixed order of coefficient derivatives by \eqref{eq:8.8}. Full slow differentiation does not
have to commute with \NSi{NS-F-P096-034}{\mathcal I_c}; Lemma~\ref{lem:8.2} estimates those
derivatives instead.
\end{EESegment}
\end{proof}

\begin{EESegment}{EE-TGT-P096-009}
\subsection{A five-dimensional correction that preserves the moment constraints.}\label{sec:8.6}
At this point, there are inverses for the pressure in Proposition~\ref{prop:8.3}, the
auxiliary-averaged tangential residual in Corollary~\ref{cor:8.5}, and the part with zero auxiliary
average in Lemma~\ref{lem:8.6}. What remains are the three integral defects
\NSi{NS-F-P096-035}{\mathcal P,\mathcal J_\theta,\mathcal J_z} from Equations~\eqref{eq:8.12}
and \eqref{eq:8.15}. This subsection corrects their linear changes by adding a velocity increment
that also leaves the two constraints in \eqref{eq:8.2} unchanged. Three azimuthal coefficients and
two axial coefficients solve the five equations in \eqref{eq:8.25}.
\end{EESegment}
